\documentclass[12pt]{article}

\usepackage[utf8]{inputenc}
\usepackage{amsmath, amssymb, amsthm}
\usepackage{geometry}
\geometry{a4paper, margin=2.5cm}

\title{Teorema de Caracterizare a Probabilităților}
\author{}
\date{}

\theoremstyle{definition}
\newtheorem{definition}{Definiție}

\theoremstyle{plain}
\newtheorem{theorem}{Teoremă}

\begin{document}

\maketitle

\begin{theorem}[Teorema de caracterizare a probabilităților]
Fie $(\Omega, \mathcal{F})$ un spațiu măsurabil și fie 
$\mu : \mathcal{F} \to [0,1]$ o funcție finit aditivă, adică


\[
    \mu(A \cup B) = \mu(A) + \mu(B)
    \quad \text{pentru } A \cap B = \varnothing.
\]


Atunci $\mu$ este o probabilitate (adică este $\sigma$-aditivă) dacă și numai dacă este continuă la creștere, adică pentru orice șir de mulțimi


\[
    A_1 \subseteq A_2 \subseteq \cdots,
\]


avem


\[
    \mu(A_n) \uparrow \mu\!\left( \bigcup_{n=1}^{\infty} A_n \right).
\]


\end{theorem}

\noindent
Echivalent, $\mu$ este probabilitate dacă și numai dacă este continuă la scădere pe șiruri cu măsură finită:


\[
    A_1 \supseteq A_2 \supseteq \cdots, \quad \mu(A_1) < \infty
    \;\Rightarrow\;
    \mu(A_n) \downarrow \mu\!\left( \bigcap_{n=1}^{\infty} A_n \right).
\]



\end{document}
